\chapter{Présentation de l'entreprise et cahier des charges}

\section*{Introduction}
\addcontentsline{toc}{section}{Introduction}

Ce chapitre vise à décrire en détail le contexte du stage de fin d'études, réalisé au sein de l'entreprise TIM Group Tunisie. Dans un premier temps, nous présentons un aperçu général de l'entreprise, en soulignant ses principaux atouts, ses missions et quelques données clés. Ensuite, nous examinons plus en profondeur les services et produits offerts par TIM Group Tunisie, ainsi que ses secteurs d'activité et ses principales références clients. Enfin, nous présentons le contexte et les objectifs du stage, suivis de la description de notre projet ainsi que des résultats attendus.

% ---------------------------------------------------------------
\section{Présentation de l'entreprise d'accueil}
% ---------------------------------------------------------------

TIM Tunisie, ou Techniques Industrielles et Management, est un organisme multidisciplinaire fondé en 2003 et spécialisé dans le domaine du management et de l'ingénierie. L'entreprise se distingue par le développement d'applications informatiques dédiées, offrant une large gamme de produits logiciels évolutifs, notamment la gestion de la qualité assistée par ordinateur (GQAO), la gestion de la maintenance assistée par ordinateur (GMAO) et la gestion de la production assistée par ordinateur (GPAO).

L'entreprise compte une dizaine de collaborateurs, en majorité des ingénieurs dans les domaines de l'informatique, et des consultants en génie industriel. TIM travaille étroitement avec des partenaires locaux et à l'international, notamment en France, au Canada et au Maroc.

\begin{figure}[H]
    \centering
    \includegraphics[height=3cm]{fig/avatar.png}
    \caption{Logo de TIM Group Tunisie}
    \label{fig:tim_logo}
\end{figure}

% ---------------------------------------------------------------
\subsection{Les atouts et les missions de l'entreprise}
% ---------------------------------------------------------------

TIM Tunisie offre une expertise spécialisée dans plusieurs domaines de management, notamment la gestion de la qualité, la santé et sécurité au travail, la protection de l'environnement, la gestion des denrées alimentaires, la production et la maintenance. L'entreprise se distingue par sa flexibilité et son agilité, son accès aux meilleures pratiques et une vision externe objective.

Les systèmes de management mis en place assurent l'efficacité attendue, la performance opérationnelle et la certification. Les services comprennent le conseil, l'accompagnement, l'audit et la formation. Ils sont fortement consolidés par la fourniture d'applications informatiques dédiées à la qualité, la production, la maintenance, et autres domaines connexes.

Les atouts majeurs de TIM résident dans sa capacité à offrir des projets intégrés combinant l'expertise métier avec les systèmes d'information appropriés. Pour ce faire, elle se base sur des ressources de haut niveau, consolidées si nécessaire par de la sous-traitance internationale spécialisée.

% ---------------------------------------------------------------
\subsection{Chiffres clés et pays d'interventions}
% ---------------------------------------------------------------

Le tableau de bord général de TIM se distingue par des indicateurs de performance significatifs :

\begin{itemize}
    \item Plus de 1\,000 projets réalisés entre consulting, formations et logiciels.
    \item Interventions dans plus de 40 secteurs, dont principalement l'agroalimentaire, l'automobile et l'oil \& gas.
    \item Une expérience de plus de 20 ans.
    \item Plus de 200\,000 personnes touchées.
\end{itemize}

TIM Tunisie intervient dans plusieurs pays, y compris la Tunisie, la Guinée Conakry, la France, le Maroc, l'Algérie, l'Italie, l'Allemagne, l'Espagne, la Bulgarie, l'Afrique du Sud, la République du Congo et la Libye.

% ---------------------------------------------------------------
\subsection{Les services, produits et références}
% ---------------------------------------------------------------

TIM Tunisie offre une gamme de services couvrant trois axes principaux :

\begin{itemize}
    \item \textbf{Consulting :} Expertise spécialisée pour établir et optimiser les systèmes de management.
    \item \textbf{Digitalisation :} Modernisation des processus opérationnels avec des solutions logicielles spécialisées.
    \item \textbf{Formation :} Programmes de formation couvrant divers domaines de management.
\end{itemize}

Les produits comprennent des solutions numériques pour la gestion de la qualité, la maintenance et la production, disponibles en mode On-Premise ou SaaS dans le Cloud. Ils incluent également des applications mobiles offrant diverses fonctionnalités accessibles depuis un smartphone ou une tablette.

\vspace{0.5em}
\noindent\textbf{QALITAS}

QALITAS est le produit phare de TIM Tunisie. Il est conçu pour optimiser la gestion des systèmes de management QHSE. Cet outil permet la digitalisation des processus de la qualité, de la santé-sécurité et de l'environnement. Il facilite ainsi la conformité aux normes internationales telles qu'ISO 9001, ISO 45001 et ISO 14001 dans leurs versions en vigueur. QALITAS offre des fonctionnalités avancées pour la gestion des audits, le suivi des indicateurs de performance (KPI), la gestion des non-conformités (NC), ainsi que la gestion des risques et opportunités (R\&O). Il propose également une interface conviviale pour une utilisation simplifiée et une intégration fluide avec d'autres systèmes informatiques.

\vspace{0.5em}
\noindent\textbf{GMAO PRO}

GMAO PRO est un logiciel de gestion de la maintenance assistée par ordinateur. Cette solution est destinée à la gestion de la maintenance des équipements industriels. Elle permet de gérer à la fois les pannes et le planning préventif des équipements, en passant par les stocks de pièces de rechange.

% ---------------------------------------------------------------
\subsection{Références de TIM}
% ---------------------------------------------------------------

TIM Tunisie collabore avec des entreprises issues de divers secteurs, notamment la cosmétique, la chimie, l'oil \& gas, la mécanique, l'électricité, l'électronique, le plastique, l'enseignement, ainsi que les domaines de l'inspection, du contrôle, de la certification, de l'agroalimentaire et du textile.

% ---------------------------------------------------------------
\section{Contexte du projet}
% ---------------------------------------------------------------

TIM Tunisie commercialise QALITAS, une plateforme de management QHSE utilisée par un nombre croissant d'entreprises en Tunisie et à l'international. QALITAS centralise la gestion des risques et des opportunités (R\&O) conformément aux exigences des normes ISO 9001, ISO 45001 et ISO 14001. Les utilisateurs y saisissent et évaluent leurs risques, définissent des actions préventives ou correctives, et suivent les indicateurs de performance via des tableaux de bord.

Cependant, TIM a constaté que la majorité des utilisateurs de QALITAS peinent à exploiter pleinement les données disponibles sur la plateforme. Le monitoring continu des risques, la réévaluation dynamique en fonction des signaux faibles, et la détection proactive des opportunités restent des activités essentiellement manuelles, chronophages et sujettes à des oublis.

Face à l'émergence des technologies d'intelligence artificielle, et en particulier des grands modèles de langage (LLM) et des systèmes multi-agents (SMA), TIM a souhaité explorer la conception d'une couche d'intelligence artificielle capable d'assister, voire d'automatiser, le cycle de pilotage des risques et opportunités au sein de QALITAS. C'est dans ce cadre que s'inscrit le présent projet de fin d'études.

% ---------------------------------------------------------------
\section{Objectifs du projet}
% ---------------------------------------------------------------

L'objectif principal du projet est de concevoir et de développer une \textbf{plateforme d'agents IA} capable d'extraire, d'analyser et d'exploiter les données issues des tableaux de bord qualité de QALITAS, afin d'assurer le monitoring continu, la détection, la réévaluation et le pilotage des risques et des opportunités.

Plus précisément, le projet vise à atteindre les objectifs suivants :

\begin{itemize}
    \item \textbf{Extraction et structuration des données :} Lire et parser automatiquement les données des dashboards QALITAS (risques, opportunités, indicateurs, non-conformités, actions) au format Excel ou via l'API REST de QALITAS.

    \item \textbf{Analyse et évaluation intelligente :} Calculer et réévaluer le score de criticité (RPN) de chaque risque en intégrant les dimensions de probabilité, impact et détectabilité, et identifier les signaux d'alerte justifiant une révision.

    \item \textbf{Détection des opportunités :} Identifier les processus performants et les conditions favorables pouvant être transformées en opportunités d'amélioration, conformément aux exigences ISO.

    \item \textbf{Génération d'actions correctives et préventives :} Proposer des plans d'actions contextualisés, priorisés par criticité et alignés sur les bonnes pratiques QSE, puis les injecter automatiquement dans QALITAS.

    \item \textbf{Pilotage continu et reporting :} Produire des rapports de synthèse structurés à destination des responsables qualité, résumant l'état des risques, les alertes actives et les actions recommandées.
\end{itemize}

% ---------------------------------------------------------------
\section{Cahier des charges}
% ---------------------------------------------------------------

Le projet se concentre sur la conception et le développement d'une architecture multi-agents intelligente composée de quatre agents spécialisés et coordonnés par un orchestrateur central.

\vspace{0.5em}
\noindent\textbf{Agent 1 -- Extraction et structuration des données}

Le premier agent est responsable de la collecte et de la structuration des données brutes issues des sources disponibles : exports Excel des tableaux de bord QALITAS, fichiers CSV de KPI, et données extraites via l'API REST de QALITAS. Il produit un contexte de données unifié exploitable par les agents suivants.

\vspace{0.5em}
\noindent\textbf{Agent 2 -- Analyse et réévaluation des risques}

Le deuxième agent analyse les risques enregistrés dans QALITAS. Il recalcule le score RPN (Risk Priority Number) de chaque risque en tenant compte de l'évolution des indicateurs de performance, de l'historique des non-conformités et des tendances observées. Il identifie les risques critiques devant faire l'objet d'une réévaluation prioritaire et ceux nécessitant la mise en place ou la mise à jour d'actions.

\vspace{0.5em}
\noindent\textbf{Agent 3 -- Détection des opportunités}

Le troisième agent analyse les données de performance des processus pour identifier des opportunités d'amélioration. Il s'appuie sur les KPI, les tendances favorables et les bonnes pratiques sectorielles pour proposer des opportunités pertinentes, argumentées et injectables dans QALITAS.

\vspace{0.5em}
\noindent\textbf{Agent 4 -- Génération d'alertes et d'actions}

Le quatrième agent génère des alertes contextualisées pour les risques critiques et propose des actions correctives ou préventives associées. Ces actions sont directement injectées dans QALITAS via l'API REST, avec les liens vers les entrées risques concernées, les responsables suggérés et les délais de mise en oeuvre.

\vspace{0.5em}
\noindent\textbf{Orchestrateur et exigences transverses}

Un orchestrateur central, implémenté avec le framework LangGraph, coordonne l'enchaînement des agents et gère le contexte partagé. Les exigences transverses incluent : la robustesse face aux données manquantes ou incohérentes, la traçabilité des décisions de l'IA par journalisation, la gestion des doublons via un cache d'injection, et la conformité avec les modèles de données de l'API QALITAS.

Le système devra être opérationnel sur un environnement Windows avec accès à la plateforme QALITAS et aux modèles de langage (LLM) locaux via Ollama ou distants via les API d'Anthropic et OpenAI. Des livrables complets incluant le code source, la documentation technique et un rapport de validation sur données réelles seront fournis à TIM.

% ---------------------------------------------------------------
\section*{Conclusion}
\addcontentsline{toc}{section}{Conclusion}
% ---------------------------------------------------------------

Au terme de ce chapitre, nous avons présenté l'entreprise d'accueil TIM Group Tunisie, ses activités, ses produits, ses services et ses principales références. Nous avons mis en évidence le contexte qui a motivé ce projet, à savoir le besoin d'une couche d'intelligence artificielle capable d'enrichir la plateforme QALITAS pour un pilotage dynamique et automatisé des risques et des opportunités. Le cahier des charges définit une architecture multi-agents claire, avec quatre agents spécialisés coordonnés par un orchestrateur LangGraph.

Dans les chapitres suivants, nous développons en détail la conception, l'implémentation et la validation de chacun de ces agents, en appuyant nos choix sur des résultats concrets obtenus sur des données réelles issues de QALITAS.
